\documentclass[12pt]{jlreq}
\usepackage{amsmath, amssymb}

\newcommand{\C}{\mathbb{C}}
\newcommand{\R}{\mathbb{R}}
\newcommand{\Z}{\mathbb{Z}}
\newcommand{\N}{\mathbb{N}}
\newcommand{\F}{\mathbb{F}}
\newcommand{\Prob}{\mathbb{P}}
\newcommand{\Exp}{\mathbb{E}}
\newcommand{\dd}{\,d}
\newcommand{\Ker}{\operatorname{Ker}}
\newcommand{\Impart}{\operatorname{Im}}
\newcommand{\Hom}{\operatorname{Hom}}

\begin{document}

\(K = \C(S, T, U)\) を複素数体 \(\C\) 上の \(3\) 変数有理関数体とし，\(L\) を \(K\) 上の多項式
\[
  F(x) = x^6 - Sx^4 + Tx^2 - U
\]
の最小分解体とする．

(1) 拡大次数 \([L : K]\) を求めよ．

(2) \(L\) に含まれる \(K\) の \(6\) 次ガロワ拡大をすべて求め，各拡大についてその拡大を生成する元を \(1\) つ与えよ．

(3) \(X \in L\) を \(F(x)\) の根とし，\(\sigma\) をガロワ群 \(G = \operatorname{Gal}(L/K)\) の位数 \(3\) の元として，\(Y = \sigma(X)\)，\(Z = \sigma(Y)\) とおく．\(M = K((X-Y)(Y-Z)(Z-X))\) を \(K\) 上 \((X-Y)(Y-Z)(Z-X)\) で生成される \(L\) の部分体としたとき，拡大次数 \([M : K]\) を求めよ．また，\(M\) が \(K\) 上のガロワ拡大であるかどうか判定せよ．

\end{document}